% Volume IV: Law as an Epistemic Machine
% Part VIII. Legal Friction as Freedom
% Chapter 42: Warrants and the Cost of Looking
% Spine: proof-spines/ch042-spine.md

\chapter{Warrants and the Cost of Looking}
\label{ch:ch042}

A warrant makes intrusion costly on purpose. Before officials may look, they must specify what they expect to find and why, converting generalized suspicion into a contestable claim that another institutional actor can inspect. The chapter represents a warrant as a \textbf{bounded search authorization}
\[
W=(T,S,Q,\Delta t),
\]
\Definition\ where \(T\) is the target, \(S\) the scope, \(Q\) the evidentiary justification, and \(\Delta t\) the temporal duration.
Each coordinate matters: the target prevents rummaging, the scope prevents drift, the justification ties looking to reasons stated in advance, and the time bound prevents an open-ended license. A warrant is thus the price law places on observation so that successful discovery cannot retroactively justify the search that produced it. The same protection turns into concealment when demands for ex ante specificity are applied so rigidly that recurring, publicly legible patterns of hidden abuse never become sufficient to authorize any scrutiny at all.
